\section{A compact base carrying a product decomposition}\label{sec:spreading}

The single-fiber decomposition theorem does not, by itself, give a
decomposition of a family.  We obtain one after a proper surjective base
change by parameterizing finite covers, factor submanifolds, and graphs.
The auxiliary base may have larger dimension than \(S\).  Its compactness,
together with a fixed compact K\"ahler ambient for the factors, will be
the essential feature.

\subsection{A marked family of klt strata}

We retain the good open \(U\), the log-smooth family, and the integer \(m\)
of the preceding sections.  In particular, the line
\(\mathcal G_m\) consists of relative log plurivolumes.  The integer \(k\)
is the largest number of horizontal coefficient-one components that meet
over \(U\); when there are none, \(k=0\).
All opens in this section are analytic Zariski opens.

\begin{lemma}\label{spread:marked}
There are a smooth compact K\"ahler manifold \(\Sigma\), a proper
surjection \(\rho:\Sigma\to S\), and a dense open
\(\Sigma^\circ\subseteq\rho^{-1}(U)\) with the following data:
\begin{enumerate}[label=\textup{(\roman*)}]
\item a proper smooth family
\(a:\mathscr A^\circ\to\Sigma^\circ\) with connected fibers and a
holomorphic section;
\item an effective relative SNC rational boundary \(D_{\mathscr A}\)
with coefficients in \([0,1)\);
\item an identification
\[
 a_*\cO_{\mathscr A^\circ}
       \bigl(m(K_{\mathscr A^\circ/\Sigma^\circ}+D_{\mathscr A})\bigr)
 \simeq \rho^*\mathcal G_m.
\]
The relative log-pluri bundle itself is the pullback of this line.
Moreover, \(\mathscr A^\circ\) and its boundary have extensions in a
smooth compact K\"ahler manifold mapping to \(\Sigma\).
\end{enumerate}
Each fiber pair is a connected component of a deepest coefficient-one
stratum of the original good family, with its adjunction boundary.
\end{lemma}

\begin{proof}
Suppose first that \(k>0\).  Choose an irreducible component \(A\) of an
intersection of \(k\) coefficient-one components whose image contains
\(U\), after decreasing \(U\) if necessary.  The original SNC hypothesis
makes \(A\) a smooth compact K\"ahler manifold.  Maximality of \(k\) says
that no further coefficient-one component meets \(A\) over \(U\).
Adjunction there gives
\[
 (K_Y+\Delta)|_A=K_A+D_A,
\]
where \(D_A\) is effective and has SNC support and coefficients less than
one.  Indeed, the remaining boundary equations are transverse coordinate
equations on \(A\).  Two distinct such equations cannot define the same
divisor on \(A\), by the SNC codimension condition.  The preparation of
\(U\) makes \(A_U\to U\), and all its boundary strata, smooth.
For \(k=0\), set \(A=Y\) and use the original boundary over \(U\).

Let \(A\to B\to X\) be the Stein factorization.  Over \(U\), the finite
map \(B\to X\) is \'etale, and \(A\to B\) has connected smooth fibers.
Resolve the graph of the map \(A\dashrightarrow S\), taking the
resolution to be unchanged over \(A_U\); call the resulting manifold
\(\Sigma\).  It is compact K\"ahler and maps holomorphically both to
\(A\) and to \(S\).  The first map also gives \(\Sigma\to B\).  Consider
\[
 A\times_B\Sigma\longrightarrow\Sigma.
\]
Over \(\Sigma^\circ=A_U\) this is the desired family: its fiber is the
connected component selected by the point of \(B\), and the map
\[
 \sigma\longmapsto \bigl(\sigma,\sigma\bigr)
\]
under the identification with \(A_U\) gives its diagonal section.
Resolve its main component, without changing this open family.
The fiber product is a closed analytic subspace of \(A\times\Sigma\),
so a projective resolution is compact K\"ahler.  Pullbacks and strict
transforms of the original boundary give analytic extensions of all
boundary data.

Ordered pluriresidue gives the indicated pullback of
\(\mathcal G_m\).  On every fiber it is a nowhere-vanishing section of
the \(m\)-th log canonical bundle, interpreted with its boundary twist.
That bundle is therefore trivial on each connected fiber, and its
space of sections has dimension one.  Grauert's base-change theorem
identifies its direct image with a line bundle.  The evaluation map is
an isomorphism because its restriction to each fiber is an isomorphism.
This proves the assertion about the actual relative bundle as well.
\end{proof}

For a section \(s\) of \(\rho^*\mathcal G_m\), denote its adjunction
plurivolume on the selected fiber by \(s_A\).  We use the norm
\begin{equation}\label{spread:residue-norm}
 \|s\|_{\mathrm{res}}
   =\left(\int_{A_\sigma}|s_A|^{2/m}\right)^{m/2}.
\end{equation}
The measure in this formula is the measure of a meromorphic
pluricanonical form with the specified boundary poles.
Since the boundary is SNC and every coefficient is less than one,
the integral is finite, and it is positive for \(s\ne0\).
The same definition will be used after any further base change.

\begin{lemma}\label{spread:residue-growth}
Let \(P\) be a smooth compact complex manifold and \(r:P\to S\) a
holomorphic map factoring through the marked-family construction
on a dense open \(P^\circ\).
Equip \(r^*\mathcal G_m\) there with
\eqref{spread:residue-norm}.  Along a disk in \(P\) mapping generically
to \(P^\circ\), generators of any sufficiently divisible positive
power of \(r^*(mM_S)\), and the dual generators, have at most
powers-of-logarithm growth in the corresponding norms, after a
sufficiently divisible power substitution.
\end{lemma}

\begin{proof}
This is Corollary~\ref{orders:residuenorm}; we recall the norm
identification.  By Theorem~\ref{orders:trace}, the corrected generator
is a power of a generator of the canonical extension of the pure Hodge line.
Proposition~\ref{hodge:residue} maps this line isomorphically to the
top residue eigenline over the chosen connected stratum component,
including all sheets over that component.  This identification
preserves canonical generators after unipotentization.

The residue is a holomorphic top form on the resolved cyclic root
cover of the klt stratum.  Its Hodge norm is its volume integral.
Its \(m\)-th power is the pullback of \(s_A\), so change of variables
identifies its squared norm with
\(\int_{A_\sigma}|s_A|^{2/m}\), multiplied by the fixed degree of the
root cover and a fixed normalization constant.  Resolution exceptional
sets have measure zero.  The canonical highest-line norm estimates
therefore give the assertion and its dual.  They apply to the
pulled-back pure line on the stated disk; no extension of a geometric
product over the center of that disk is required.
\end{proof}

\subsection{Splitting a fiber and realizing its finite covers}

We first specify the single-fiber consequence that will be spread.

\begin{lemma}\label{spread:fiber-splitting}
Let \((F,D)\) be a smooth connected compact K\"ahler pair, with
effective SNC rational boundary of coefficients less than one, and
suppose that \(m(K_F+D)\) is trivial.  There is a finite \'etale cover
\(\nu:J\to F\) and an isomorphism of pairs
\begin{equation}\label{spread:fiber-product}
 (J,\nu^*D)\simeq (Q,C)\times T_1\times\cdots\times T_j,
\end{equation}
where \((Q,C)\) is a smooth projective klt SNC pair, and each \(T_i\)
is a positive-dimensional complex torus or an irreducible holomorphic
symplectic manifold.  The boundary is entirely on \(Q\).
A point is allowed for \(Q\), and \(j=0\) is allowed.  Moreover,
\(m(K_Q+C)\) and all \(mK_{T_i}\) are trivial.
\end{lemma}

\begin{proof}
Linear triviality implies numerical triviality, so the decomposition
theorem of Matsumura--Wang--Wu--Zhang applies to this smooth effective
klt pair
\cite[version~1, Corollary~1.3]{MWWWZ25}.
It gives a product of a rationally connected pair and boundary-free
torus, Calabi--Yau, and holomorphic symplectic factors after a finite
\'etale cover.  The covered manifold and all its product factors are
smooth.  Apply the ordinary
Beauville--Bogomolov decomposition to any boundary-free factors still
requiring decomposition, and compose the further finite \'etale
covers \cite[Theorem~1]{Beauville83}.

For completeness, a rationally connected compact K\"ahler manifold
\(R\) is projective.  To see the relevant vanishing directly, graphs
in the Douady space of \(\PP^1\times R\) parameterize its holomorphic
rational curves.  There are countably many finite-dimensional
parameter families \cite{Fujiki79}.  Matching marked endpoints
parameterizes chains, again in countably many analytic families.
Stratify their parameter spaces into smooth pieces.  General pairs
of points of \(R\) occur as endpoints, so Sard's theorem implies that
one two-endpoint evaluation
\[
 e=(e_0,e_\infty):T\longrightarrow R\times R
\]
is submersive at some point.  Otherwise all these countably many
images would have measure zero.  Mark the ends of each chain segment
by \(0,\infty\), using reparameterization of \(\PP^1\).
For \(\alpha\in H^0(R,\Omega_R^2)\), its pullback by each segment
\(\PP^1\times T\to R\) comes from \(T\), because \(\PP^1\) has no
holomorphic forms of positive degree.  Comparing successive endpoints
gives
\[
 e^*(\operatorname{pr}_1^*\alpha-\operatorname{pr}_2^*\alpha)=0.
\]
Submersivity forces \(\alpha=0\) on an open set, hence everywhere.
Thus \(H^{2,0}(R)=0\).  Hodge decomposition now makes every real
degree-two class of type \((1,1)\).  Openness of the K\"ahler cone
gives a rational K\"ahler class, and an integral multiple is the first
Chern class of a positive line bundle.  Kodaira's embedding theorem
makes \(R\) projective \cite{Kodaira54}.

The same criterion makes every remaining strict Calabi--Yau factor
with \(H^{2,0}=0\) projective.  Dimension-two symplectic factors are
kept among the \(T_i\).  Combine all projective factors into \(Q\).
The boundary on \(Q\) remains effective klt SNC under this product
description.  Finally, restrict the trivial log-pluri bundle of
\(J\) to factor slices.  The canonical lines of the complementary
factors contribute only constant one-dimensional vector spaces,
so these restrictions give \(m(K_Q+C)\sim0\) and \(mK_{T_i}\sim0\).
\end{proof}

\begin{lemma}\label{spread:finite-covers}
Let \(g:E^\circ\to B^\circ\) be a proper smooth holomorphic map of
connected complex manifolds, with connected compact fibers and a
section.  Suppose \(B^\circ\) is a dense open in a smooth compact
K\"ahler manifold \(B\), and the family has a smooth compact
K\"ahler extension \(E\to B\).
Then all finite connected \'etale covers of individual fibers occur
in countably many families, each obtained after a finite \'etale base
cover of \(B^\circ\).  Each resulting base and total space admit
smooth compact K\"ahler extensions preserving the open families.
The same holds with pulled-back relative SNC boundary data.
\end{lemma}

\begin{proof}
Fix a base point and write \(\Pi=\pi_1(E_b)\) and
\(\Gamma=\pi_1(B^\circ)\), with fiber base points supplied by the
section.  Ehresmann's theorem makes \(g\) a differentiable fiber
bundle.  Its homotopy sequence contains
\[
 \pi_2(E^\circ)\longrightarrow\pi_2(B^\circ)
 \longrightarrow\Pi\longrightarrow\pi_1(E^\circ)
 \longrightarrow\Gamma\longrightarrow1.
\]
The section makes the first arrow surjective, so the connecting
homomorphism into \(\Pi\) is zero.  It also splits the last
homomorphism.  Thus there is a genuine semidirect product
\begin{equation}\label{spread:semidirect}
 \pi_1(E^\circ)\simeq\Pi\rtimes\Gamma.
\end{equation}

The group \(\Pi\) is finitely generated, since the fiber is a compact
manifold.  For each integer \(e>0\), it has only finitely many
subgroups of index \(e\): their coset actions are represented among
the homomorphisms from a fixed finite generating set to the finite
group of permutations of \(e\) elements.  If \(H\subseteq\Pi\) is
one such subgroup, its stabilizer \(\Gamma_H\) under the action in
\eqref{spread:semidirect} has finite index in \(\Gamma\).
Pass to the corresponding finite base cover.  The subgroup
\[
 H\rtimes\Gamma_H\subseteq\Pi\rtimes\Gamma_H
\]
then defines a finite connected cover of the pulled-back total
space.  Its restriction to a fiber is the connected cover
corresponding to \(H\).  Topological covering charts lift the complex
structure, so this is a holomorphic \'etale cover.  Transport to
other fibers, and then all indices \(e\), realizes every required
fiber cover among countably many such constructions.

Extend the finite base cover first across \(B\setminus B^\circ\)
by the analytic finite-cover extension theorem, and resolve the
normal extension.  Pull back \(E\), resolve its main component, and
extend the finite total-space cover over that resolved compact
space in the same way.  Projective resolutions give smooth compact
K\"ahler manifolds, unchanged on the open under consideration
\cite[Lemma~2.21]{Villadsen24}.

One can also describe the extension locally after making the
complement SNC.  A punctured coordinate chart has commuting
meridians.  Sufficiently divisible powers of its boundary coordinates
kill their finite permutation action.  The cover becomes a disjoint
union of trivial covers on the punctured power chart; the finite
normal extensions descend and glue by uniqueness.  Thus this step
uses finite-cover extension, for which essential singularities do
not arise.  Finite maps over K\"ahler targets and projective
resolutions preserve the K\"ahler property; equivalently one may
use the compactification in the cited Lemma.  Boundary pullbacks
and their strict transforms give the required analytic closures.
\end{proof}

\subsection{Compact parameter spaces}

We call a subset of a compact complex space \emph{constructible} if it
is a finite union of differences of closed analytic subsets.
Compactness is important in the following elementary form of analytic
constructibility.

\begin{lemma}\label{spread:constructible}
Let \(q:Z\to T\) be a proper holomorphic map of compact complex
spaces.  The image of a constructible subset of \(Z\) is
constructible in \(T\).
Consequently, any finite combination of incidence, emptiness,
surjectivity, and injectivity conditions on fibers of compact
analytic diagrams is constructible.
\end{lemma}

\begin{proof}
It suffices to consider \(A\setminus C\), where \(C\subseteq A\)
are closed analytic subsets of \(Z\).  There are finitely many
irreducible components of \(A\), so we may assume that \(A\) is
irreducible and \(C\subsetneq A\).  Set \(T'=q(A)\), a compact
irreducible analytic subspace by the proper mapping theorem.
On a dense open of \(T'\), the fibers of \(A\to T'\) have dimension
\(\dim A-\dim T'\).  Each component of \(C\) either has proper image
in \(T'\), or has smaller generic fiber dimension.  The proper
fiber-dimension theorem therefore gives a dense open
\(T'_0\subseteq T'\) on which no fiber of \(A\to T'\) is contained
in \(C\).  Hence
\[
 T'_0\subseteq q(A\setminus C).
\]
The part of the image over \(T'\setminus T'_0\) is the image of
\[
 \bigl(A\cap q^{-1}(T'\setminus T'_0)\bigr)\setminus C
\]
under a proper map whose image has smaller dimension.  Induction
proves that this remainder is constructible.  The case of a
zero-dimensional image is immediate.  This proves the first
assertion.

For example, failure of containment of two fiber subspaces is the
image of their difference.  Failure of surjectivity is detected by
points of the target outside the image, itself analytic under a
proper map.  Failure of injectivity is detected on the fiber product
of the source with itself, after removing its diagonal.  Every
ambient projection here is proper; the first assertion applies
to the indicated differences.  Finite Boolean combinations remain
constructible.
\end{proof}

We also use smoothness and rank conditions in these diagrams.
A universal Douady or Hilbert family is proper and flat.
Its nonsmooth locus is analytic, and its proper image gives the
complement of the smooth-fiber locus.  On this locus, fiber dimensions
are locally constant.  Ranks of maps of relative tangent bundles
can be tested using relative differentials and their Fitting ideals.
Thus they too give constructible conditions, also after restricting
to a constructible incidence locus.  A smooth compact fiber is
connected exactly when \(h^0(\cO)=1\), which is an analytic
semicontinuity condition in a proper flat family.  For marked
divisors, smoothness and the ranks of their defining differentials
test the relative SNC condition, including the expected codimensions
of all intersections.

\begin{proposition}\label{spread:product}
After a proper surjection \(r:P\to S\), with \(P\) a smooth compact
connected manifold, there is a dense open \(P^\circ\subseteq r^{-1}(U)\)
on which the marked klt family has a finite \'etale cover and a product
decomposition
\begin{equation}\label{spread:family-product}
 (\mathscr J,D_{\mathscr J})
 \simeq
 (\mathscr Q,C)\times_{P^\circ}
 \mathscr T_1\times_{P^\circ}\cdots\times_{P^\circ}\mathscr T_j .
\end{equation}
Here all families are proper and smooth with connected fibers.
They have the following properties.
\begin{enumerate}[label=\textup{(\roman*)}]
\item The pair \((\mathscr Q,C)\) is a projective family of effective
klt SNC pairs.  It is embedded in \(P^\circ\times\PP^N\) and is the
pullback of universal marked families by the restriction of a
holomorphic map from \(P\) to a fixed product of projective Hilbert
schemes.
\item Each \(\mathscr T_i\) is a family of submanifolds of a fixed
compact K\"ahler manifold \(V\), of positive dimension \(n_i\).
Its cohomological local systems have their integral lattices modulo
torsion and flat real polarizations.  Its top Hodge line has rank one
and is given by
\[
 \lambda_i
 =(\pi_i)_*\omega_{\mathscr T_i/P^\circ}
 =F^{n_i}R^{n_i}(\pi_i)_*\C
\]
Each index is of one of the following two types:
\[
 \begin{array}{ll}
 \text{torus type:}&
 h^{1,0}(T_i)=n_i,\quad
 \bigwedge^{n_i}F^1R^1(\pi_i)_*\C
       \xrightarrow{\ \sim\ }\lambda_i;\\[2mm]
 \text{symplectic type:}&
 n_i\ \text{even},\quad h^{2,0}(T_i)=1,\quad
 \bigl(F^2R^2(\pi_i)_*\C\bigr)^{\otimes n_i/2}
       \xrightarrow{\ \sim\ }\lambda_i .
 \end{array}
\]
The displayed maps are cup products on the underlying variations.
\item On every fiber, \(m(K_Q+C)\) and \(mK_{T_i}\) are trivial.
With
\[
 \mathcal G_{m,Q}
   =(\pi_Q)_*\cO_{\mathscr Q}
                  \bigl(m(K_{\mathscr Q/P^\circ}+C)\bigr),
\]
there is an identification of actual holomorphic line bundles
\begin{equation}\label{spread:volume-line}
 r^*\mathcal G_m
   \simeq \mathcal G_{m,Q}\otimes
                   \bigotimes_i\lambda_i^{\otimes m}.
\end{equation}
\item If the finite cover in \eqref{spread:family-product} has degree
\(e\), the norm in \eqref{spread:residue-norm} satisfies
\begin{equation}\label{spread:norm-product}
 \|s\|_{\mathrm{res}}
   =c\,\|s_Q\|_{Q,\mathrm{res}}\,
                   \prod_i\|\omega_i\|_{\mathrm{Hdg}}^m
 \quad\text{when}\quad
 s\longleftrightarrow
       s_Q\otimes\bigotimes_i\omega_i^{\otimes m}.
\end{equation}
Here \(c>0\) is constant; with compatible volume conventions,
\(c=e^{-m/2}\).  The norm on the left gives generators and their duals
of sufficiently divisible powers of \(r^*(mM_S)\) the two-sided
logarithmic bounds of Lemma~\ref{spread:residue-growth}.
\end{enumerate}
These assertions are preserved on the good open under further smooth
compact auxiliary base changes.  Neither generic finiteness of \(r\)
nor a global K\"ahler form on \(P\) is required.
\end{proposition}

\begin{proof}
Apply Lemma~\ref{spread:finite-covers} to the marked family from
Lemma~\ref{spread:marked}.  We obtain countably many compact K\"ahler
total spaces \(V_\alpha\to B_\alpha\), whose good fibers supply all
finite \'etale fiber covers that may occur in
Lemma~\ref{spread:fiber-splitting}.  Each \(B_\alpha\) maps properly
to \(S\).  Fix one such family temporarily, and write \(V\to B\).
The boundary on its good open is the restriction of a fixed rational
divisor on \(V\); only its restriction to good fibers will be tested.

\smallskip
\noindent\emph{The parameters and their geometric conditions.}
For an individual splitting, embed \(Q\) into some \(\PP^N\).
Choose points in complementary factors and realize each \(T_i\)
as a factor slice of \(J\subset V\).  The product isomorphism is
represented by its graph in
\[
 V\times\PP^N\times V^j.
\]
We parameterize the fiber \(J\) by \(B\), the embedded \(Q\) and its
marked boundary components by Hilbert schemes of \(\PP^N\), the
slices by Douady spaces of \(V\), and the graph by the Douady space
of the displayed ambient.

For each choice of \(N,j\), dimensions, boundary coefficients,
Hilbert polynomials, and Douady components, take the product
\[
 K=B\times\mathscr H\times
       \prod_{i=1}^j\mathscr D_i\times\mathscr D_\Gamma,
\]
where \(\mathscr H\) is the corresponding product of Hilbert schemes.
This is a compact complex space.  Indeed, connected components of
the Douady space of a compact K\"ahler manifold are compact
\cite[Corollary~5.3]{Toma11}, and all ambients here are compact
K\"ahler.  Fujiki's countability theorem \cite{Fujiki79} gives
countably many components.  There are also only countably many
choices of the discrete data and of \(\alpha\).  Thus all the
individual splittings have been included in countably many compact
parameter spaces.

In each \(K\), impose the following conditions on the universal data:
\begin{enumerate}[label=\textup{(\alph*)}]
\item the base parameter is good; the projective factor, the slices,
and the graph are nonempty smooth fibers of the prescribed
dimensions; the factors are connected;
\item every marked boundary component is a smooth divisor on the
projective factor, distinct marks have distinct supports, and all
their intersections have the prescribed SNC codimensions;
\item the slices lie in \(V_b\), and the graph lies in the product
of \(J=V_b\) with the indicated projective factor and slices;
\item the projections of the graph to \(J\) and to the factor
product are bijective, with invertible vertical differentials;
\item the graph identifies the weighted boundary on \(J\) with the
boundary pulled back from the marked projective factor.
\end{enumerate}
These conditions define a constructible subset of \(K\).
For (a)--(b), use proper flat smoothness, connectivity, and the
relative rank tests described above.  Condition (c) is incidence.
For (d), all ambient coordinate projections are holomorphic on
the compact universal graph.  Bijectivity is tested by
Lemma~\ref{spread:constructible}, and the vertical rank conditions
are imposed by relative differentials.  For (e), group marks of
equal coefficient and test equality of the corresponding reduced
supports on the graph.  The boundary extensions on \(V\) are
analytic, so these are again incidence conditions on compact
diagrams.  On the SNC locus, equality by coefficient groups is
exactly equality of the rational divisors.

Decompose this constructible locus into finitely many irreducible
locally closed analytic pieces and then stratify their singular loci.
This yields countably many smooth connected locally closed analytic
strata, each with compact analytic closure.  Over every such stratum,
the graphs give isomorphisms of smooth families: fiberwise they are
bijective and have invertible vertical differential, so the relative
inverse function theorem and properness give holomorphic inverses.

\smallskip
\noindent\emph{The cohomological types and domination.}
Retain the strata containing at least one splitting with a specified
list of torus and symplectic types.  This is still a countable
collection; no constructibility of a Hodge locus is being assumed.
We check that the cohomological conditions in (ii) then hold
throughout the chosen stratum.

For a smooth factor family over such a stratum \(T\), pull back a
K\"ahler form \(\omega_V\) from \(V\).  It is relatively K\"ahler,
and its fiberwise class is flat for the Gauss--Manin connection,
because it is represented by a closed form on the total family.
The smooth proper K\"ahler Hodge theorem therefore supplies the
usual pure variations and locally constant Hodge numbers; one may
apply \cite[Theorem~1.1(i)]{FF25} with empty removed boundary on
small smooth base charts, and use purity of compact K\"ahler cohomology.
Fiberwise Lefschetz decomposition for this flat class is a
decomposition of real variations.  Poincar\'e duality on the
primitive summands, with the Hodge--Riemann signs, gives flat real
polarizations on the full cohomology.  This construction requires
no rational K\"ahler class.  Locally over \(T\), adding a K\"ahler
form from a coordinate neighborhood to the ambient form makes
the total space K\"ahler as well.

The cup products
\[
 \bigwedge^{n_i}R^1(\pi_i)_*\C
      \longrightarrow R^{n_i}(\pi_i)_*\C,
 \qquad
 \bigl(R^2(\pi_i)_*\C\bigr)^{\otimes n_i/2}
      \longrightarrow R^{n_i}(\pi_i)_*\C
\]
are flat morphisms of pure variations of equal weight in each case.
Their kernels and images are subvariations; equivalently,
polarizations give flat Hodge complements and strictness gives
the induced Hodge filtrations.  The ranks on the highest Hodge
pieces are consequently locally constant.  At a torus fiber,
exterior multiplication of holomorphic one-forms gives its
nonzero top form.  At an irreducible holomorphic symplectic fiber,
the top power of the symplectic form gives its nonzero top form.
Thus the Hodge numbers and cup-product isomorphisms in (ii)
hold throughout \(T\).  Only these cohomological consequences of
the two types are needed.

Every point of the original good base has splitting data on at
least one of the retained strata: first select the connected
stratum fiber, then its finite cover, then its product and
embeddings.  The compact closures of these strata have proper
analytic images in \(S\).  If all their images were proper analytic
subsets, their countable union could not contain the nonempty open
\(U\), by Baire category.  One closure therefore dominates \(S\).
Resolve that irreducible closure, choosing a projective resolution
unchanged on its smooth open stratum.  The result is a smooth
compact connected \(P\), with a proper surjection \(r:P\to S\).
Let \(P^\circ\) be the inverse image of the chosen stratum.
The product of families on it is
\eqref{spread:family-product}.  Projection to \(\mathscr H\) gives
the holomorphic Hilbert-scheme map on all of \(P\) required by (i).
The ambient \(V\) is fixed by the compact parameter space containing
this closure.  This proves (i)--(ii).

\smallskip
\noindent\emph{Volumes and their norms.}
The finite cover of the selected klt fiber preserves the trivial
\(m\)-th log canonical bundle.  Restriction of
\eqref{spread:family-product} to factor slices proves the fiberwise
trivialities in (iii), even for parameter points where only the
cohomological types have been retained.

Grauert's base-change theorem gives the direct image line
\(\mathcal G_{m,Q}\).  For each \(T_i\), the Hodge-number calculation
gives \(\rk\lambda_i=1\).  A nonzero holomorphic top form on \(T_i\)
is nowhere vanishing: its \(m\)-th power is a nonzero section of a
trivial line bundle on a connected compact manifold, hence has
no zeros.  Consequently, relative multiplication induces an
isomorphism
\[
 \lambda_i^{\otimes m}
       \simeq(\pi_i)_*\omega_{\mathscr T_i/P^\circ}^{\otimes m}.
\]
Finite \'etale pullback of the selected adjunction volume, followed
by the product formula for relative canonical bundles and their
boundary twists, now gives \eqref{spread:volume-line}.  Every map
in this construction identifies actual one-dimensional spaces
of sections on the fibers, so these are holomorphic line-bundle
isomorphisms, with their specified identification on the open.

Take a decomposable local section on the right side of
\eqref{spread:volume-line}.  With compatible volume conventions,
change of variables under the \'etale cover and Fubini's theorem
give
\[
 e\int_{A_\sigma}|s_A|^{2/m}
   =
 \left(\int_{Q_\sigma}|s_Q|^{2/m}\right)
       \prod_i\left(\int_{T_{i,\sigma}}|\omega_i|^2\right).
\]
The integral on \(Q_\sigma\) is finite by the klt SNC condition.
Raising the equality to the power \(m/2\) proves
\eqref{spread:norm-product}.  The degree \(e\) is fixed on the
connected open parameter stratum.  Other fixed volume conventions
only change the constant \(c\).
Lemma~\ref{spread:residue-growth} gives the assertion for
\(r^*(mM_S)\) and its dual.

All constructions are compatible with pullback on their smooth
opens.  The argument has used compactness of \(P\) for proper
images and later global arguments, and a K\"ahler form on \(V\)
for the relative Hodge theory.  It has imposed no dimension
restriction on \(r\) and no global K\"ahler hypothesis on \(P\).
Empty products and zero-dimensional projective factors satisfy
the same formulas, with integration over a point taken with mass one.
\end{proof}
